\section{Four numbers, and what happened to them}
\label{sec:worked}

\noindent\emph{Status: draft -- a ledger, not a result. The four quantities below have all been quoted in this
report at some time, and they are now sorted by whether they can still be quoted. Two of them cannot, one only
after being renamed, and one can with a caveat attached to its reference impedance. The values are printed here
exactly as the records give them; the sorting is the content. The rules applied are those of
Section~\ref{sec:method}.}

A measurement that is withdrawn is more useful than a measurement that is quietly repaired. When a number stops
being defensible, the temptation is to adjust it -- to change the normalisation, to choose a longer window, to
report the value from the run that agrees -- and to keep the sentence it supported. Doing so destroys the only
evidence that the earlier reasoning was wrong, and it hides the cost. This section therefore keeps every value
and changes only the status word beside it: a number may be quotable, diagnostic, or withdrawn, and the value is
never the thing that is corrected. In two of the four cases below the status moved a long way while the digits
did not move at all, and that is exactly the outcome this ledger exists to display.

\subsection{The ledger}
\label{sec:worked-ledger}

\begin{table}[H]
  \centering
  \caption{Four quantities from this report's records, with provenance and the status they now carry.
  \emph{Quotable}: the record supports the quantity as a measurement of what its name says. \emph{Diagnostic
  only}: the quantity is exact and repeatable, but it measures a property of the sampling rather than of the
  device, and is reported under a name that says so. \emph{Withdrawn}: the quantity was quoted as a measurement
  and is not one; the value is retained for the audit trail and must not be cited. The provenance column is
  part of the number: a value whose run, mesh and estimator are not stated here is not reproducible.}
  \label{tab:worked}
  \small
  \begin{tabular}{@{}p{3.1cm}p{3.2cm}p{6.2cm}p{2.0cm}@{}}
    \toprule
    Quantity & Value & Provenance: run, mesh, estimator & Status \\
    \midrule
    (a) Transmission coefficient, phase constant and reference impedance of the reference box at
    \SI{5.8278}{\giga\hertz}
    & $s_{21} = 0.98088 + 0.00476\,j$;\newline $\gamma = 0.01022 + 0.00418\,j$;\newline
      $z_0 = 49.76 + 13.59\,j$\,\si{\ohm};\newline $\beta = 229.554$\,m$^{-1}$
    & CW run, mesh $132\times104\times38$ with a \SI{0.2}{mm} cell; driven at \SI{5.827786408327}{\giga\hertz}
      against a requested \SI{5.827864}{\giga\hertz}; $600$ of $2400$ steps discarded, $1800$ accumulated over
      $6.8636695303115883\times10^{-10}$\,s; two-wave fit at a transverse cut
    & quotable \\
    \addlinespace
    (b) Station-profile tilt of the time-averaged Poynting flux
    & $0.987038$ as a ratio;\newline $6.48$ per cent as a tilt
    & CW run at \SI{5.8278}{\giga\hertz}; nine stations at cells $52$--$100$ of the propagation axis, \SI{9.6}{mm}
      span; flux formed by time-averaging $\mathbf{E}\times\mathbf{H}$ at each station, then
      $r = \sqrt{s_9/s_1}$
    & diagnostic only \\
    \addlinespace
    (c) ``Field loss'' of the reference box
    & $16$ per cent, later ``at most'' $2$ per cent
    & the same nine-station arrays as (b), read through two different normalisations; no separate run exists
    & withdrawn \\
    \addlinespace
    (d) Passivity deficit from a paired $|S_{11}|$ and $|S_{21}|$
    & $7.33$ per cent
    & its provenance is not established in this report
    & withdrawn, pending \\
    \bottomrule
  \end{tabular}
\end{table}

Three of the four rows are the same record, and that is the point of the table rather than a defect in it. One
field dump was processed four ways: a two-wave fit at a cut, a nine-station flux profile, a loss figure derived
from that profile, and -- in the fourth row -- an unrelated published pair that turned out not to reproduce. The
first is quotable, the second is a diagnostic, the third is withdrawn, and the fourth is withdrawn pending. The
digits in rows (a) and (b) have not changed since they were first written down. Only the claims attached to them
have, and the next three subsections say why for each.

A status is assigned by three tests, and they are mechanical enough to be checked by a reader. First, does the
derivation of the quantity assume the proposition the quantity is supposed to establish? A flux ratio that
assumes a smooth envelope fails this test whatever its value. Second, is the quantity a ratio of two numbers
obtained by the same operation on the same field dump? If so it may be stable to many digits and still say
nothing about level, and it cannot be quotable on its own. Third, does the record carry a second, independent
route to the same quantity? Only where all three tests pass does the word quotable apply, and the tests are
applied to the quantity as it is used in a sentence, not to the quantity as it might in principle be computed.
Row (b) is an exact number and fails the first test; rows (c) and (d) fail it outright; row (a) passes all
three, which is why it is the only row that is quotable.

\subsection{Quotable: one controlled transmission record}
\label{sec:worked-cw}

The record of row (a) is the best-controlled run in the report, and it is quotable because every assumption in
its reduction was checked against something else in the same record. The mesh is $132\times104\times38$ cells
with a \SI{0.2}{mm} cell, which puts the transverse cross-section at roughly a tenth of a guide wavelength; the
drive was snapped to the nearest frequency the discrete time step can represent on the transform grid, which is
\SI{5.827786408327}{\giga\hertz} against a requested \SI{5.827864}{\giga\hertz}, an offset of
$-77\,591.67$\,Hz, or $-13.3$\,ppm; the step is $\dd t = 3.8131497390619936\times10^{-13}$\,s; and the first
$600$ of $2400$ steps were discarded, leaving $1800$ accumulated steps spanning $6.8636695303115883\times10^{-10}$\,s.
The offset of $-13.3$\,ppm is quoted rather than absorbed because it is a property of the drive: a resonance
measured $-13.3$\,ppm away from the requested frequency is the resonance of a slightly different structure,
and the reader is entitled to know by how much.

At that drive the two-wave fit at a transverse cut returns
\begin{equation}
  \begin{aligned}
    s_{21} &= 0.9808843848072781 + 0.004761855437949522\,j , \\
    \gamma &= 0.010215720318428398 + 0.004178365747705708\,j ,
  \end{aligned}
  \label{eq:worked-s21}
\end{equation}
with $|\gamma| = 0.0110371953841$ and a fit residual of $0.0128478440321$. The reflection coefficient is the
reason this record is quotable and the others are not: at $|\gamma| \approx 0.011$ the standing-wave ratio on
the line is $1.022$, so the field at a cut is very nearly a single travelling wave, and the two-wave model is
forced by the record rather than imposed on it. In the records of Section~\ref{sec:msres-alpha} the same model
returns a rate that disagrees with the joule audit by a factor of $1.86$; here there is no fitted rate at all,
only a complex amplitude at a cut, and the quantity being quoted is that amplitude.

Two numbers in the same reduction need their status stated separately, because both look like measurements and
only one is. The effective permittivity is $\eps_{\mathrm{eff}} = 3.53220159465$, and it is a restatement of the
phase constant: by $\eps_{\mathrm{eff}} \equiv (\beta c/\omega)^2$, once $\beta = 229.554353095$\,m$^{-1}$ is
known at a known $\omega$ the permittivity is arithmetic, and it cannot disagree with the phase constant by so
much as a digit. It is recorded because a reader wants it and because the dispersion curves of
Section~\ref{sec:dispersion} are drawn against it; it is not a second measurement of anything, and any sentence
of the form ``the phase constant and the effective permittivity agree'' would be a sentence about a
substitution.

The reference impedance is the number with the caveat. The fit returns
\begin{equation}
  z_0 = 49.76097429242066 + 13.592389198926165\,j \,\si{\ohm} ,
  \qquad |z_0| = \SI{51.5839859517}{\ohm} \text{ at } 15.277852^{\circ} ,
  \label{eq:worked-z0}
\end{equation}
and the imaginary part is not noise. An imaginary term of \SI{13.6}{\ohm} in the reference impedance says that
the line cannot be referred to a real \SI{50}{\ohm} reference without a transformation: the reflection
coefficient, and therefore every $S$-parameter computed from this record, depends on which reference impedance
is assumed, and ``the transmission coefficient of this line'' is not a well-defined number until the reference is
named. That is not a defect of the fit and it is not a large effect -- the angle is $15^{\circ}$, so
$\operatorname{Re} z_0$ is within a per cent of \SI{50}{\ohm} -- but it is the kind of caveat that has to travel
with the number rather than be rediscovered by whoever uses it next. Accordingly the quotable claim is the pair
of amplitudes of equation~\eqref{eq:worked-s21} together with the reference of
equation~\eqref{eq:worked-z0}, and any statement of an attenuation or a mismatch derived from them is only as
good as the reference it was derived against.

The record also carries its own conservation check, and it is this rather than the smallness of $|\gamma|$ that
lets the third test of Section~\ref{sec:worked-ledger} be passed. From the amplitudes of
 equation~\eqref{eq:worked-s21},
\begin{equation}
  |s_{21}|^{2} + |\gamma|^{2} \;=\; 0.96228 ,
  \label{eq:worked-cw-passivity}
\end{equation}
so $3.77$ per cent of the incident power is neither returned nor transmitted and is accounted for as absorbed in
the section. That is not a second measurement of the transmission coefficient -- it is the same two amplitudes
entered once more -- but it is an independent check on the \emph{pair}, because passivity is a property of the
structure and not of the fit: a fit that returned a transmission above unity, or a reflection inconsistent with
the standing-wave ratio at the cut, would fail here even if each amplitude separately looked plausible. It is
also the number that the pairing of Section~\ref{sec:worked-pairing} fails to reproduce, which is why the two
rows are worth reading side by side: one record, two amplitudes, and a deficit of $3.77$ per cent on one line of
arithmetic, against a published $7.33$ per cent that cannot be recovered from it.

\subsection{Diagnostic only: the station-profile tilt}
\label{sec:worked-tilt}

The second number is the one this section exists for, because it is the only number in the report that was
demoted without any digit changing, and because the demotion was forced by an identity rather than by a
disagreement.

The construction is this. The field is sampled at nine stations spaced six cells apart along the propagation
axis, \SI{1.2}{mm} apart, spanning \SI{9.6}{mm}. At each station $k$ the time-averaged Poynting flux $s_k$ is
formed by averaging $\mathbf{E}\times\mathbf{H}$ over the accumulated window, and the flux at the last station is
divided by the flux at the first to give a transmission-like ratio $r = \sqrt{s_9/s_1}$. If the flux declined
only by absorption, then $r = |S_{21}|$ and the record measures attenuation. The reason it does not is that the
two quantities $r$ and the fractional drop of the flux are related by an identity in which the attenuation does
not appear at all.

Let $\tau$ be the tilt of the station profile, the drop from first to last measured against the mean of the
array, and $\bar{s}$ that mean. Then
\begin{equation}
  1 - r^{2} \;=\; \frac{s_1 - s_9}{s_1} \;=\; \tau \cdot \frac{\bar{s}}{s_1} ,
  \qquad
  \tau \;=\; \frac{s_1 - s_9}{\bar{s}} .
  \label{eq:worked-tilt}
\end{equation}
Nothing in equation~\eqref{eq:worked-tilt} refers to a material, a cross-section or a rate. It is an algebraic
relation between three averages of the same nine numbers, and it says that whatever $r$ reports, it reports a
property of the shape of the station profile, scaled by the ratio of the mean to the first value. If the profile
declines linearly, $\bar{s}/s_1$ is close to one and $1 - r^2$ is nearly the tilt; if the profile is a standing
wave sampled at an arbitrary phase, the same is still true, and the number that comes out is the amplitude of
the standing wave, not the loss.

The CW record at \SI{5.8278}{\giga\hertz} gives, at the nine stations lying on cells
$52$, $58$, $64$, $70$, $76$, $82$, $88$, $94$ and $100$,
\begin{equation*}
  \begin{aligned}
    s_1, \dots, s_9 \;=\; {} & 6614.038753,\ 6584.568766,\ 6563.085584, \\
    & 6543.527531,\ 6516.123028,\ 6470.735308, \\
    & 6400.792767,\ 6306.335991,\ 6194.943464 ,
  \end{aligned}
\end{equation*}
with a mean of $6466.0168$. The last station over the first is $s_9/s_1 = 6194.943464/6614.038753 = 0.936635$, a
power ratio, so the power deficit is $6.3365$ per cent, printed as $6.34$. The amplitude transmission is its
square root, $r = \sqrt{s_9/s_1} = 0.96780$, and a deficit taken on that basis is $3.22$ per cent. On the same
array the tilt is $6.4815$ per cent and the mean-to-first factor is $0.977617$, and their product,
$6.4815 \times 0.977617 = 6.3364$ per cent, against the measured $6.3365$ per cent: a residue of $0.0001$
per cent, the size of rounding the tilt to four decimal places, so the product reproduces the power deficit
at the quoted precision and not to the last printed place. That agreement is not a coincidence to be
explained; it is equation~\eqref{eq:worked-tilt} holding, as it must, because the three numbers are three
averages of one array. The reason the deficit and the
tilt are so close is that the profile is nearly straight over nine points, which is what a fraction of a
standing-wave period looks like.

A deficit quoted as a bare percentage is therefore ambiguous until it says whether the ratio it came from was a
power ratio or an amplitude ratio. Here the two differ by close to a factor of two: $6.34$ per cent on the power
ratio $0.936635$, $3.22$ per cent on the amplitude ratio $0.96780$. That is one way a number acquires two
defensible-looking values, and it is why this report prints the ratio as well as the percentage.

The same extraction on the pulse record at \SI{5.333}{\giga\hertz} gives a deficit of $1.9561$ per cent against a
tilt of $1.9757$ per cent and a mean-to-first factor of $0.9901$; the product is $1.9561$ per cent, and there the
identity closes to every printed digit. On the substrate-only record the nine stations are
\begin{equation*}
  \begin{aligned}
    s_1, \dots, s_9 \;=\; {} & 4879.172642,\ 4867.110209,\ 4862.878825, \\
    & 4859.305586,\ 4847.939993,\ 4820.684887, \\
    & 4772.478055,\ 4703.721293,\ 4621.547925 ,
  \end{aligned}
\end{equation*}
with a mean of $4803.871$, a ratio of $0.97324$ ($5.27$ per cent), a tilt of $5.363$ per cent and a
mean-to-first factor of $0.984568$, whose product is $5.28$ per cent. Three records, three arrays, and in every
one the ``attenuation'' is the tilt multiplied by a factor within a few per cent of one.

That is the whole of the diagnostic, and the geometry says why it is a diagnostic. The stations span
\SI{9.6}{mm}. At \SI{5.8278}{\giga\hertz} the free-space wavelength is \SI{51.449}{mm}, and with the
$\eps_{\mathrm{eff}} = 3.5322$ of the controlled record of Section~\ref{sec:worked-cw} the guide wavelength of
the cross-section they sample is $\lambda_g = \lambda_0/\sqrt{\eps_{\mathrm{eff}}} =
\SI{51.449}{mm}/1.879415 = \SI{27.375}{mm}$, so the span is $0.35\,\lambda_g$. The flux on a line carrying a
forward and a reflected wave ripples at twice the phase constant, that is, with period
$\lambda_g/2 = \SI{13.688}{mm}$ in space, so the array covers seventy per cent of one period of that ripple. The
first and last stations of such a sample can fall at \emph{any} relative phase, including the phase that makes
them differ by several per cent while the line is lossless, and seven tenths of a period is nowhere near enough
of the ripple to average it out. A reading that does not average the standing wave is confounded at any loss
level, including zero.

What that reading was called is the history this subsection records. The quantity was for a period written as
$|S_{21}|_{\text{poynting}} = 0.987038$ and quoted as an attenuation of $-0.1133$\,dB over the \SI{9.60}{mm}
span, which is the largest single contribution the report ever claimed towards the loss of the reference box,
and it was presented as model-free -- which, in the sense of ``containing no fitted parameter'', it was. It is
now called the station-profile tilt, and it carries \emph{no} attenuation claim. The value $0.987038$ survives
in the table and in this sentence; the name does not, because the name is what the arithmetic does not support.
The correspondence is exact: $20\log_{10}(0.987038) = -0.1133$\,dB, so nothing has been recomputed, only
relabelled, which is the honest form of a correction.

Two further consequences follow from the relabelling, and both were visible only afterwards. The first is that
the tilt cannot serve as a bound on the loss either. At resonance, the full-section tilt of $6.48$ per cent is
larger than the circuit-level passivity deficit of $3.77$ per cent computed from the amplitudes of
equation~\eqref{eq:worked-s21}; a quantity that exceeds the thing it is supposed to bound is not a bound, it is
a different quantity. On the substrate-only record the tilt is $5.36$ per cent against the same $3.77$ per cent,
and there it does bound the deficit -- which is the one surviving quantitative role of the construction, and the
reason it is a diagnostic worth keeping rather than deleting. The second consequence is that the tilt is a good
measurement of something else, namely of how much standing wave is present at the array: a tilt far above the
passivity deficit is evidence that the sampling sits inside a ripple, which is exactly the warning one wants
before quoting any between-station rate. The array was doing useful work all along. It was just not doing the
work it was named for.

\subsection{Withdrawn: the field-loss figures}
\label{sec:worked-withdrawn}

Two published figures are withdrawn outright: a ``$16$ per cent field loss'' of the reference box, and, before
that, a figure of ``at most $2$ per cent''. Both were computed from the nine-station arrays of
Section~\ref{sec:worked-tilt}, through two different normalisations of the same flux profile, and both are
therefore readings of the tilt rather than of a loss.

The withdrawal is not a small correction and it is worth saying plainly what it removes. Between ``at most $2$
per cent'' and ``$16$ per cent'' lies a factor of eight, and both numbers were at some point the report's answer
to the question of how much the reference box loses. They differ because the two normalisations respond
differently to where the standing wave sits in the array -- which is the signature of a sampling artefact and not
of physics, since a physical loss cannot change by a factor of eight when nothing about the geometry changes.
The audit of Section~\ref{sec:worked-tilt} explains both figures with one mechanism, arithmetic on nine
numbers, and that explanation is the reason they are withdrawn rather than demoted. A diagnostic is a quantity
that is exactly what it is; these two were never exactly anything, because the normalisation was chosen after
the fact to make the number plausible.

What replaces them is not another reading of the same array. The rate of loss that the report now quotes for the
reference box comes from the joule audit of Section~\ref{sec:ms-accounting}, in which the power dissipated in
the conductors is integrated over the volume of the metal and divided by the power entering the section, and it
is $0.01160$\,dB/mm. That number is not a ratio of two samples of the same quantity and it does not depend on
where the standing wave sits; it is the work the walls actually did. The station arrays remain, and their tilt
remains, as the audit trail for the withdrawal -- and as a warning attached to any future flux-based reading:
nine stations spanning a third of a guide wavelength cannot measure a loss, whatever they are called.

\subsection{Withdrawn pending: the passivity pairing}
\label{sec:worked-pairing}

The fourth row is different from the other three in every respect: it comes from a different reduction, it does
not involve the station arrays, and it is withdrawn pending a question that has not been answered rather than
withdrawn because the answer is known.

A published pairing of $|S_{11}|$ and $|S_{21}|$ at the reference box recorded a passivity deficit of $7.33$ per
cent, that is, $|S_{11}|^2 + |S_{21}|^2$ short of unity by that fraction at resonance. A single CW record now
carries both coefficients at the same frequency, and it gives
\begin{equation}
  |s_{21}|^{2} + |\gamma|^{2} \;=\; 0.96228 ,
  \label{eq:worked-passivity}
\end{equation}
a deficit of $3.77$ per cent. The two deficits do not agree: they differ by a factor close to two, far outside
anything that the per-cent-level uncertainties in either record could explain. The row therefore does not
reproduce, and its status is withdrawn pending -- withdrawn because the number it quotes cannot be defended,
pending because the run that produced it has not been found and the question of what it actually measured is
still open.

The correction to the record concerns not the number but the \emph{reason} I first gave for the mismatch, and it
is recorded here because the reason was quoted and was wrong. My first note on the discrepancy asserted that the
pairing \emph{could not} have come from a single run -- that no one field dump could yield both coefficients at
the resonance, so the pairing must be a mixture of two records. That claim is too strong and it is false. A
single CW run does yield both numbers at $f_0$: $\gamma$ from the two-wave fit at a cut, $s_{21}$ from the same
fit, both referred to the same plane. What the single record shows is not that the pairing is impossible but
that it is different: $3.77$ per cent against $7.33$ per cent.

The difference between the two kinds of failure is the reason the correction matters. A provenance
impossibility would be closed by finding the missing run, and the number would then be either vindicated or
replaced. A numeric mismatch is closed only by re-measuring, and until that is done the honest entry is an open
discrepancy with a stated value. By asserting impossibility I had moved the failure into the category with the
easier remedy, and made the report look more rigorous while making it less true -- the second error described in
Section~\ref{sec:method-provenance}, committed in the same direction as the first. What stands now is the weak
statement: the pairing does not reproduce against the single record, the reason is unknown, and the row stays
withdrawn until it is either located or replaced.

\subsection{What the ledger says}
\label{sec:worked-lessons}

One row of four is quotable without qualification. That is the honest count for this stage of the work, and it is
worth reading against the length of the report: the great majority of the numbers in this document are
restatements, diagnostics or arithmetic, and the four entries above are the ones that were load-bearing enough
to audit.

The audit did not leave the argument where it found it, and this section should not be read as claiming that
nothing was lost. Three things were. The ``two model-free readings agree'' finding of
Section~\ref{sec:ms-criterion} is no longer a finding, because one of the two was the tilt; the coincidence that
the tilt lands $1.9$ per cent from the audited rate is recorded and not counted, and the criterion is now
satisfied on the joule audit alone. The $0.671$ fall of the ratio in Section~\ref{sec:msres-verdict} is quoted
and not read, since the quantity that falls is the tilt. And the claim about where the missing loss goes rests
on the audited work per unit length rather than on a comparison of two readings. Each of these is a weaker
statement than the one it replaced, and each is one that a reader can check.

What the ledger buys is the opposite of what it costs. Rows (a) and (b) are the same record, and after the audit
they can both be used -- one as a measurement, one as a diagnostic -- without either contaminating the other,
because the border between them is written down. The alternative, keeping the tilt under its old name, would have
preserved a paragraph of agreement and would have made every future comparison between the report's loss figures
and the literature an unmarked comparison of unlike things. A number that is defendable is worth a paragraph of
argument, and the four rows of Table~\ref{tab:worked} are the accounting.
